\scanpage{21}
as a quasi-projective scheme over $S$.

2. Theorem of the cube. I believe we discussed about it time ago, but maybe the proof I told you was valid only if one assumes the Pic functor of one of the \tikz[baseline=(s.base)]{\node[inner sep=0pt](s){\deleted{\unsure{functors}}};\node[anchor=south west,inner sep=1pt] at ([xshift=0em,yshift=6pt]s.north west){\inserted{factors}};} involved representable. To prove unramified-
ness of the functor \uline{Corr} however you need only a week infinit\rlap{i}es\rlap{i}emal form of the theorem of the square, for which you will find a proof in the manuscript notes I am joining on correspondance classes, containing also the proof\deleted{\unsure{s}} of the statements you were recalling in your question 4. I hope you will be able to read them, I agree the handwriting is wretched and the notes moreover very sketchy. - On the other hand, I recall you that the theorem of the cube follows rather formally once one knows separatedness of $\underline{\mathrm{Corr}}_{S}(X,Y)$ for two of the three factors involved, and using the usual formal properties of the Picard functor (among which commutation with inverse limits of Artin rings is the less trivial).

3. As for the separatedness of $\underline{\mathrm{Corr}}_{S}(X,Y)$, this is about trivial whene-
ver the Pic functor of one of the factors $X,Y$ is separated? Now this is \tikz[baseline=(s.base)]{\node[inner sep=0pt](s){certainly};\draw[gray,line width=.35pt] ([xshift=-3pt,yshift=2pt]s.south west)--([xshift=-2.4pt,yshift=-2pt]s.south west);\draw[line width=.55pt] ([xshift=-3pt,yshift=2pt]s.south west)--([xshift=-2.8pt,yshift=.5pt]s.south west);} the case \deleted{if} for $X$ if its geometric fibers are integral, (a fortiori if $X$ is an abelian scheme over $S$ !). To show this, one may assume $S$ the spectrum of a \tikz[baseline=(s.base)]{\node[inner sep=0pt](s){valuation};\node[anchor=south west,inner sep=1pt] at ([xshift=-1em,yshift=6pt]s.north west){\inserted{discrete}};} ring, and one is \rlap{d}reduced to show that if $\underline{L}$ is an invertible sheaf on $X$ whose restriction to \tikz[baseline=(s.base)]{\node[inner sep=0pt](s){the};\fill ([xshift=5pt,yshift=2pt]s.north east) circle (.3pt);\fill ([xshift=7pt,yshift=5pt]s.north east) circle (.3pt);} general fiber $X_{t}$ is trivial, then $\underline{L}$ is trivial. Now $X_{1}$ is an open subset of $X$, and the assumption on $\underline{L}$ can be expressed by saying that $\underline{\mathord{\rlap{$X$}L}}$ is defined by a Cartier divisor whose support is contained in the special fiber $X_{o}$. Now $X_{o}$ itself is already a Cartier divisor (d\rlap{i}efined by a global equation $\mathord{\rlap{$f$}t}=0$) and moreover is an integral subscheme of $X$, from this follows that the divisor $D$ is a multiple of $X_{o}$ (assume for
% archive sheet number: 20
